% Generated draft: requires mathematical and source review.
\documentclass[11pt]{article}
\usepackage[margin=25mm]{geometry}
\usepackage{fontspec}
\setmainfont{Latin Modern Roman}
\usepackage{amsmath,amssymb,mathtools,unicode-math}
\setmathfont{Latin Modern Math}
\usepackage{enumitem,needspace,xurl,hyperref}
\setlength{\parindent}{0pt}
\setlength{\parskip}{5pt}
\setlength{\emergencystretch}{3em}
\newcommand{\N}{\mathbb{N}}
\newcommand{\Z}{\mathbb{Z}}
\newcommand{\Q}{\mathbb{Q}}
\newcommand{\R}{\mathbb{R}}
\newcommand{\C}{\mathbb{C}}
\newcommand{\F}{\mathbb{F}}
\newcommand{\A}{\mathbb{A}}
\newcommand{\PP}{\mathbb{P}}
\newcommand{\E}{\mathbb{E}}
\newcommand{\eps}{\varepsilon}
\newcommand{\dd}{\,\mathrm{d}}
\begin{document}
\section*{UnsolvedMath ID: 20002340}
\textbf{Source problem code:} catalogue entry
\subsection{Definitions and mathematical statement}
Fix a finite field $\F_q$, and put $R=\F_q[W,Z]$ and $S=\F_q[Z]$. Question 3 (Demeyer) asks whether there are a fixed finite integer $n$ and a Diophantine predicate $\alpha(f,\bar g)$ over $R$, with $f\in R$ and $\bar g\in S^n$, such that
\[
\forall f\in R\ \exists\bar g\in S^n\ \alpha(f,\bar g),\qquad
\forall\bar g\in S^n\quad |\{f\in R:\alpha(f,\bar g)\}|<\infty.
\]
A Diophantine predicate is an existential finite system of polynomial equations over the ambient ring, with its coefficient and parameter convention specified. The source fixes neither a uniformity requirement across $q$ nor the exact named coefficients. The fiber bound is pointwise finiteness, not a uniform numerical bound. The source says that an analogous covering by finite fibers is possible if both occurrences of ``$\alpha$ holds'' are replaced by ``$\alpha$ does not hold''; it supplies no construction for that complementary assertion.

The source further asserts that a positive answer would imply that recursively enumerable sets over $R$ are precisely Diophantine sets (the DPRM property, in a fixed effective presentation, with arities understood). The research aid distinguishes this assertion from what its pullback argument establishes: singleton fibers would suffice, while arbitrary finite fibers require a further separation argument not supplied in the summary. Thus the claimed implication is retained as source context with that proof gap flagged, not silently weakened or independently certified.

The aid's bounded-degree result is explicit. For each fixed integer $D\geq0$, choose $N=q^k>D$ and use one parameter $g\in S$. The equations
\[
f=\sum_{i=0}^{D}W^i a_i,\qquad
g=\sum_{i=0}^{D}Z^i a_i^{N}
\]
with existential $a_0,\ldots,a_D\in R$ give a Diophantine relation whose nonempty $f$-fibers over $g$ are singletons and which covers every $f$ with $\deg_W f\leq D$, including $f=0$. Indeed the distinct residues $i$ modulo $N$ in the second equation force each $a_i\in S$ and determine it uniquely; its Frobenius exponent is fixed. The degree bound is fixed for the formula, so this is not one cover of all $R$.

Conversely, the aid asserts that for any fixed finite system
\[
P_j(f,\bar g,\bar y)=0\quad(1\leq j\leq r),\qquad
\bar g\in S^n,\quad\bar y\in S^m,
\]
with polynomial coefficients in $R$, if all $f$-fibers of the resulting existential relation are finite, then there exists a single integer $B$ such that every nonzero $f$ in the relation has $\deg_W f\leq B$. This is a restriction on predicates whose parameters and \emph{all} auxiliary witnesses lie in $S$; the ordinary AIM question permits ambient-ring witnesses, so the restriction does not refute it. The supplied account reports the full question open at the research run.
\subsection{Short English statement}
Demeyer asks for a Diophantine finite-fiber cover of the bivariate polynomial ring using tuples from its univariate subring. The aid codes each fixed degree stratum with singleton fibers and gives a degree obstruction when all auxiliary data stay univariate. It leaves the unrestricted question and the general finite-fiber-to-DPRM justification unresolved.
\subsection{Sources}
\begin{itemize}
\item Local UnsolvedMath registry, permanent ID $20002340$, \texttt{catalogue entry}; \texttt{aim-logic-notes.json}, \texttt{notes[115]}.
\item American Institute of Mathematics, \emph{Extensions of Hilbert's tenth problem}, Question 3, attributed to Demeyer; moderated by B. Poonen and T. Scanlon, notes by J. Demeyer. \url{https://aimath.org/WWN/hilberts10th/hilberts10th.pdf}. Supplied verification places the question on the first page. The extracted ``12'' and capitalization of the moderator line are page/header artifacts.
\item Research attempt 1: partial result, problem open at run, rigor label ``proved,'' novelty confidence low. Registry classification partial progress and status partially solved; statement status exact. No authoritative external citation for the partial results or the DPRM implication is supplied.
\end{itemize}
\end{document}
